\section{Introduction}
Social virtual reality (VR) applications, which allow multiple people to interact in virtual spaces, are among the most popular VR experiences. For example, VRChat regularly hosts 40,000 players per month \cite{vrchatmetrics}. However, the complex nature of social VR, where key visual information changes frequently, makes it particularly challenging for blind and low vision (BLV) people. Current social VR platforms do not have accessibility features to address these challenges.

Previous research on VR accessibility for BLV users has largely focused on adding spatial audio or haptic feedback to objects and actions, such as footstep sounds signifying avatar movement or haptic vibrations signifying collisions \cite{andrade2018echo, zhao2018cane, siu2020cane, zhang2020cane, dang2023opportunities}. For example, Dang et al. proposed a dual-track audio system that allows BLV users to hear spatial audio alongside concurrent audio descriptions of important actions during VR concerts \cite{dang2023opportunities}. However, this research exclusively targets single-user VR experiences. These solutions would likely be sensorially overwhelming and disruptive in the dynamic environments of social VR, requiring a fundamentally different, high-level approach. This highlights the need for a solution that can provide real-time, contextual information about scenes, rather than relying solely on low-level sensory feedback.

Recently, some researchers have addressed social VR accessibility using guiding interactions as a template \cite{collins2023gaze, wieland2022nvc, ji2022vrbubble, collins2023guide}. Previously, we proposed a sighted guide framework for BLV users in social VR, inspired by physical-world sighted guides, where a BLV person holds a sighted person's elbow for navigation support while the sighted person walks slightly ahead and to the side of them \cite{collins2023guide}. In this framework, a human sighted guide joins a BLV user to assist them in exploring virtual environments. We showed that virtual guides were a promising, easy-to-learn accessibility approach, allowing BLV users to successfully navigate virtual environments with a technique already familiar to them from the physical world. While highly effective, human guides suffer from major limitations: they are not always available, and BLV users may prefer using automated tools to maintain their independence \cite{collins2023guide, gonzalez2024scenedesc}. To overcome these challenges, the field requires an effective, on-demand guidance system that preserves a user’s autonomy, such as AI-driven tools.

To address challenges with human guides, we presented an AI guide for social VR in follow-up work \cite{collins2024ai}. We designed a guide that used a large language model (LLM) to answer users’ visual questions and provide navigation support inside VR. The guide also offered six “personas,” or combinations of appearances and mannerisms, to allow for a customizable guidance experience. Crucially, we have not yet evaluated this system's efficacy via a user study. This lack of empirical data leaves critical questions about BLV people's experience with an AI guide unanswered: Would BLV users be able to effectively utilize it? What behaviors would they exhibit toward the AI guide? Are AI-generated descriptions of social VR useful to BLV people in practice? Understanding the answers to these questions is necessary for designing an effective AI guide.

To address this need, we studied our previously developed AI guide’s \cite{collins2024ai} use in social VR environments with BLV users. In this work, our research questions are broad:  

\begin{itemize}
    \item \textit{RQ1: How is an AI guide utilized by BLV people in social virtual reality?}
    \item \textit{RQ2: How do BLV people behave towards an AI guide in social virtual reality?}
\end{itemize}

To answer these questions, we conducted a study with 16 BLV participants who utilized the AI guide in social VR environments modeled after parks. We selected two of the guide's personas—a dog and a robot—for them to experience. Participants completed two tasks in VR: 1) using the guide to become familiar with the layout of each park, and 2) leading park tours for confederates who joined the VR scene. This second task required participants to balance interaction with the environment, other users, and the AI guide. After all tasks were completed, we interviewed participants about their experience.

We found that participants were able to use the guide to explore the parks and share accurate information about their features with confederates. However, participants’ interactions with the guide varied significantly based on context. For example, participants treated the guide as a tool to address their visual needs when using it alone, but began role-playing with it around confederates. Many participants also seemed less satisfied with the guide’s descriptions when with confederates, rephrasing its descriptions with their own imagined details, such as invented backstories for avatars in the parks. We also observed that some participants treated the guide as a source of general knowledge, such as asking it for the correct terms of park features. Finally, we discuss how the observed use of the AI guide differs from use of human guides and identify opportunities for studying the use of AI guides over time.

We implemented the first AI guide based on feedback from and tested with BLV users. Through this, we contribute the first findings on an AI guide’s effectiveness with assisting BLV people in social VR and how they react to its assistance. This work offers a critical foundation for designing future AI-powered VR accessibility tools that are genuinely useful and empowering.
